\documentclass[10pt,a4paper]{article}
\usepackage[T1]{fontenc}
\usepackage[utf8]{inputenc}
\usepackage[czech]{babel}
\usepackage{lmodern}
\usepackage{amsmath,amssymb}
\usepackage{geometry}
\usepackage{enumitem}
\usepackage{array,tabularx,booktabs}
\usepackage{xcolor}
\usepackage{fancyhdr}
\usepackage{tikz}
\usepackage[hidelinks]{hyperref}
\geometry{left=22mm,right=22mm,top=22mm,bottom=22mm}
\definecolor{accent}{HTML}{243D65}
\definecolor{muted}{HTML}{596575}
\setlength{\parindent}{0pt}
\setlength{\parskip}{3pt}
\setlist[itemize]{topsep=2pt,itemsep=1pt,leftmargin=6mm}
\setlist[enumerate]{topsep=2pt,itemsep=2pt,leftmargin=7mm}
\pagestyle{fancy}\fancyhf{}\fancyhead[L]{\small\color{muted} Příprava na vyučovací hodinu}\fancyhead[R]{\small\color{muted} Kombinatorika}\fancyfoot[C]{\thepage}
\renewcommand{\headrulewidth}{0.3pt}
\newcommand{\cas}[1]{\textcolor{accent}{\bfseries #1}}
\newcommand{\otazka}[1]{\emph{#1}}
\begin{document}
{\LARGE\bfseries Úvod do kombinatoriky}\par
{\large První vyučovací hodina -- variace, permutace a kombinace}\par
\vspace{2mm}
\textbf{Délka:} 45 minut \quad \textbf{Forma:} řízená diskuse, společné řešení, samostatná práce

\textbf{Cíl hodiny:} Žák systematicky vypisuje možnosti, používá jednoduché pravidlo součinu a u konkrétních situací rozhodne, zda záleží na pořadí a zda jsou dovolena opakování. Intuitivně odliší variaci, permutaci a kombinaci.

\section*{Časový plán}
\begin{tabularx}{\textwidth}{@{}p{24mm}X@{}}\toprule
0--7 min & Motivace: uspořádání studentů a výběr dvojice \\
7--17 min & Systematický výpis, pravidlo součinu a PIN \\
17--30 min & Rozhodování: záleží na pořadí? Jsou povolena opakování? \\
30--39 min & Pojmenování základních typů: variace, permutace, kombinace \\
39--45 min & Samostatné úlohy, společná kontrola a shrnutí \\
\bottomrule\end{tabularx}

\section*{1. Motivace \texorpdfstring{\cas{(0--7 min)}}{(0--7 min)}}
Na tabuli napište: \otazka{Máme tři studenty: Adama, Barboru a Cyrila. Kolika různými způsoby se mohou postavit vedle sebe na společnou fotografii?}

Nechte studenty nejprve tipovat a poté vypisovat možnosti. Zapisujte systematicky:
\[
ABC,\quad ACB,\quad BAC,\quad BCA,\quad CAB,\quad CBA.
\]
Výsledek: \(6\) pořadí.

Navazující otázka: \otazka{Kolika způsoby z těchto tří studentů vybereme dvojici, která půjde pro pomůcky?}

Možnosti: \(AB, AC, BC\), tedy \(3\). Ptejte se: \otazka{Proč je první výsledek šest a druhý jen tři? Je dvojice Adam--Barbora jiná než Barbora--Adam?}

\textbf{Pointa:} Při fotografování záleží na pořadí, při výběru dvojice nezáleží.

\section*{2. Pravidlo součinu \texorpdfstring{\cas{(7--17 min)}}{(7--17 min)}}
Zadání: \otazka{Ve školní jídelně jsou na výběr tři hlavní jídla a dva nápoje. Kolik různých obědů lze sestavit, vybíráme-li právě jedno jídlo a jeden nápoj?}

Nejprve nechte vypsat kombinace, poté ukažte zkrácený postup:
\[
3\cdot 2=6.
\]
\textbf{Pravidlo součinu:} Jestliže první volbu můžeme provést \(m\) způsoby a pro každou z těchto voleb druhou volbu \(n\) způsoby, pak obě volby dohromady lze provést \(m\cdot n\) způsoby.

Doplňující otázka: \otazka{Kolik čtyřmístných PINů lze vytvořit z číslic 0--9, jestliže se mohou opakovat?}
\[
10\cdot10\cdot10\cdot10=10^4=10\,000.
\]
\textbf{Poznámka:} PIN může začínat nulou. Zdůrazněte, že opakování číslic je zde přípustné. Pokud by šlo o čtyřmístné \emph{číslo}, první cifra nulou být nemůže.

\section*{3. Pořadí a opakování \texorpdfstring{\cas{(17--30 min)}}{(17--30 min)}}
U každé situace mají studenti pouze rozhodnout (bez nutnosti počítat), zda záleží na pořadí a zda se prvky mohou opakovat. Nechte je zdůvodnit odpověď.

\begin{tabularx}{\textwidth}{@{}Xcc@{}}\toprule
\textbf{Situace} & \textbf{Pořadí?} & \textbf{Opakování?} \\\midrule
Ze šesti studentů vybíráme tříčlennou komisi. & Ne & Ne \\
Šesti běžcům udělujeme zlato, stříbro a bronz. & Ano & Ne \\
Vytváříme třímístný kód z číslic 0--9, opakování povoleno. & Ano & Ano \\
Ze 5 příchutí zmrzliny vybíráme 3 kopečky; opakování dovoleno, pořadí nerozhoduje. & Ne & Ano \\
\bottomrule\end{tabularx}

\textbf{Otázky do diskuse:} \otazka{Je výběr Adam--Barbora--Cyril stejná komise jako Cyril--Adam--Barbora? Může běžec dostat dvě medaile? Je kód 123 stejný jako 321?}

\section*{4. Základní typy uspořádání a výběrů \texorpdfstring{\cas{(30--39 min)}}{(30--39 min)}}
Teprve nyní pojmenujte známé situace. Zatím pracujte především s případy \emph{bez opakování}.

\begin{tabularx}{\textwidth}{@{}p{26mm}X X@{}}\toprule
\textbf{Typ} & \textbf{Charakteristika} & \textbf{Příklad} \\\midrule
Variace & Uspořádaný výběr \(k\) z \(n\) prvků & První tři místa v závodě \\
Permutace & Uspořádání všech \(n\) prvků & Seřazení všech knih na polici \\
Kombinace & Neuspořádaný výběr \(k\) z \(n\) prvků & Výběr členů komise \\
\bottomrule\end{tabularx}

Upozorněte, že u výběrů existují také varianty s opakováním. \textbf{V první hodině nezavádějte všechny vzorce:} cílem je situaci správně rozpoznat, nikoli mechanicky dosazovat.

\section*{5. Závěrečné ověření \texorpdfstring{\cas{(39--45 min)}}{(39--45 min)}}
Studenti řeší samostatně. V závěru krátce proberte postup.
\begin{enumerate}
\item Kolika způsoby lze ze čtyř různých knih vybrat dvě, pokud na pořadí nezáleží?
\item Kolika způsoby lze ze čtyř různých knih obsadit první a druhé místo na polici (bez opakování)?
\item Kolik čtyřmístných PINů lze vytvořit z číslic 0--9, pokud se žádná číslice nesmí opakovat?
\end{enumerate}
\textbf{Řešení pro učitele:} 1. \(\binom42=6\); 2. \(4\cdot3=12\); 3. \(10\cdot9\cdot8\cdot7=5040\).

\section*{Zápis do sešitu}
\textbf{Kombinatorika} se mimo jiné zabývá určováním počtu možností výběru a uspořádání prvků.

\textbf{Pravidlo součinu:} Pokud první volbu můžeme provést \(m\) způsoby a pro každou z nich druhou volbu \(n\) způsoby, potom dohromady existuje \(m\cdot n\) možností.

Při řešení kombinatorické úlohy se ptáme: \emph{Záleží na pořadí? Mohou se prvky opakovat? Vybíráme všechny prvky, nebo pouze některé?}

\textbf{Variace:} uspořádaný výběr části prvků. \quad
\textbf{Permutace:} uspořádání všech prvků. \quad
\textbf{Kombinace:} neuspořádaný výběr části prvků.

\section*{Metodická poznámka a návaznost}
První hodina má vytvořit představu, \emph{co} vlastně počítáme. V dalších hodinách je vhodné probrat pravidlo součtu, faktoriál a postupně permutace, variace a kombinace (nejdříve bez opakování, potom s opakováním). U slovních úloh vždy vyžadujte krátké zdůvodnění výběru postupu.
\end{document}
